\scanpage{128}
\margin{\fbox{\parbox{1.9cm}{Montrer qu'on \unsure{peut se ramener} à la dimension finie}}}
La première partie du théorème est résumée dans les
formules (1) à (4), \deleted{et} \inserted{outre} les affirmations relatives
aux topologies vectorielles des espaces $\mathcal L^p(\mathcal H)$.
\inserted{Montrons \deleted{\unsure{d'abord}} d'abord que ces dernières sont conséquences très faciles \deleted{\unsure{de|des}} \unsure{des formules} (3) et (4)\ldots}
La formule (1) est évidente, \unsure{car|La} \deleted{\unsure{on sait que}}
\deleted{\unsure{bien} \unsure{au} cas $p<+\infty$, et} \deleted{\unsure{pour $p<+\infty$}}, on sait (théorème
de Riesz) que $|A|$ et $|A^*|$ sont \deleted{\unsure{de même suite}} \inserted{unitairement équivalents}
(?) \deleted{\unsure{valeurs propres}}. La formule (2) est un cas
particulier de la formule (6), mais \deleted{\unsure{peut se|peut être}} \inserted{\unsure{on peut|Je peux}}
donner une démonstration \deleted{\unsure{pourtant directe}}
\deleted{démontrer simplement de façon directe}
directe, en procédant comme dans $[\,]$, .

En effet, si $(\unsure{k_i|\alpha_i})$ \deleted{\unsure{resp\textsuperscript{t} $(\beta_i)$}} désigne la
suite des valeurs \inserted{propres} de \deleted{$\unsure{\lvert BA\rvert|\lvert BA\rvert^2}$} $|A|$
\inserted{et $(\unsure{\beta_i|h_i})$ la suite des valeurs propres de $\unsure{\lvert BA\rvert|\lvert B\rvert A\rvert|\lvert BA\rvert^2}$ (resp\textsuperscript{t} $|AB|$)} \deleted{\unsure{resp\textsuperscript{t} $|BA|$}},
rangées par ordre de grandeur
décroissante (chacune répétée autant de fois
que le comporte sa multiplicité) on a
\[
\beta_i\leq\|B\|\alpha_i\quad\text{pour tout }i.\quad(\text{loc cité, page \quad, lemme \quad}),
\]
d'où aussitôt les formules (2).

La \deleted{\unsure{dernière} des formules (3) est évidente, la}
Les formules (3) \deleted{peut} \inserted{peuvent} être regardées \deleted{de la formule} comme
\uline{première} \deleted{\unsure{est une conséquence immédiate}}
\unsure{conséquences|conséquence} de la formule (6), (que nous démontrerons de façon indépendante) de la façon
suivante. Si $p'$ est tel que $\frac{1}{p}+\frac{1}{p'}=1$,
et si $B\in\mathcal L^{p'}(\mathcal H)$, \deleted{et tel que}
\inserted{$\|B\|_{p'}\leq1$}, alors
on a en vertu de (6) $\|AB\|_1\leq\|A\|_p$, d'où
\[
|\operatorname{Tr}(AB)|\leq\|A\|_p,\quad\text{et}\quad
\operatorname*{Sup}_{\|B\|_{p'}\leq1}|\operatorname{Tr}AB|\leq\|A\|_p
\]
